% Generated draft: requires mathematical and source review.
\documentclass[11pt]{article}
\usepackage[margin=25mm]{geometry}
\usepackage{fontspec}
\setmainfont{Latin Modern Roman}
\usepackage{amsmath,amssymb,mathtools,unicode-math}
\setmathfont{Latin Modern Math}
\usepackage{enumitem,needspace,xurl,hyperref}
\setlength{\parindent}{0pt}
\setlength{\parskip}{5pt}
\setlength{\emergencystretch}{3em}
\newcommand{\N}{\mathbb{N}}
\newcommand{\Z}{\mathbb{Z}}
\newcommand{\Q}{\mathbb{Q}}
\newcommand{\R}{\mathbb{R}}
\newcommand{\C}{\mathbb{C}}
\newcommand{\F}{\mathbb{F}}
\newcommand{\A}{\mathbb{A}}
\newcommand{\PP}{\mathbb{P}}
\newcommand{\E}{\mathbb{E}}
\newcommand{\eps}{\varepsilon}
\newcommand{\dd}{\,\mathrm{d}}
\begin{document}
\section*{UnsolvedMath ID: 701740}
\textbf{Source problem code:} catalogue entry
\subsection{Definitions and mathematical statement}
For a finite group \(K\), define its spectrum by
\[
\omega(K)=\{\operatorname{ord}(x):x\in K\}.
\]
Two finite groups are isospectral if their spectra are equal. The socle \(\operatorname{Soc}(K)\) is the subgroup generated by all minimal nontrivial normal subgroups of \(K\); set \(\operatorname{Soc}(1)=1\). A group \(L\) is solvable if its derived series, defined by \(L^{(0)}=L\) and \(L^{(n+1)}=[L^{(n)},L^{(n)}]\) for integers \(n\geq 0\), satisfies \(L^{(m)}=1\) for some integer \(m\geq 0\).

Let \(G\) be any finite group satisfying
\[
\text{for every finite group }H,\qquad
\omega(H)=\omega(G)\ \Longrightarrow\ H\cong G.
\]
Must the quotient \(G/\operatorname{Soc}(G)\) be solvable?
\subsection{Short English statement}
If a finite group is uniquely determined up to isomorphism among finite groups by its set of element orders, must its quotient by its socle be solvable?
\subsection{Sources}
\begin{itemize}
\item Local UnsolvedMath registry: permanent ID \(701740\), source problem code \texttt{catalogue entry}, title \emph{Kourovka Notebook Problem 20.82}.
\item Attributed in the supplied record to D. O. Revin, proposed in 2022.
\item Supplied reference: \emph{The Kourovka Notebook (Unsolved Problems in Group Theory)}, 20th issue (2022), Problem 20.82; also identified with the 21st issue edition (2026), \texttt{arXiv:1401.0300v48}. The collection record names E. I. Khukhro and V. D. Mazurov as editors. These bibliographic details are supplied attribution and have not been independently verified.
\item The registry labels the problem \emph{open}; this is a supplied status label, not an independently verified assessment.
\end{itemize}
\end{document}
